\documentclass[conference]{IEEEtran}
\IEEEoverridecommandlockouts
\usepackage[utf8]{inputenc}
\usepackage[english]{babel}
\usepackage{cite}
\usepackage{amsmath,amssymb,amsfonts}
\usepackage{graphicx}
\usepackage{textcomp}
\usepackage{xcolor}
\usepackage{hyperref}
\usepackage{booktabs}
\usepackage{float}

\graphicspath{{assets/}}

\begin{document}

\title{Autonomous Multi-Agent Systems for UAV Flight Safety Inspections using Large Language Models}

\author{\IEEEauthorblockN{Burak YORUK}
\IEEEauthorblockA{\textit{Department of Computer Engineering} \\
\textit{Ege University}\\
İzmir, TURKEY \\
ORCID: 0000-0001-8980-191X}
}

\maketitle

%% ============================================
%% ABSTRACT
%% ============================================
\begin{abstract}
Unmanned Aerial Vehicle (UAV) operations often rely on manual checklists and human intuition for flight safety, which introduces the risk of human error during complex or stressful missions. This paper introduces a fully autonomous Multi-Agent System (MAS) architecture designed to conduct real-time UAV flight safety inspections using Large Language Models (LLMs). By integrating live telemetry data—such as wind speed, battery health, and GPS satellite counts—with an interactive pre-flight checklist, the proposed system dynamically assesses operational risk.

Our architecture is based on the Foundation for Intelligent Physical Agents (FIPA) Agent Communication Language (ACL) standards. The system orchestrates four specialized agents: a Coordinator, a Risk Assessor, a Flight Advisor, and a Report Writer. The Risk Assessor autonomously queries environmental telemetry via APIs and continuously evaluates whether conditions satisfy safety constraints, issuing preemptive NO-GO commands even if the human operator manually verifies a checklist item. We demonstrate the practical feasibility of this approach through a lightweight web-based interface built with Flask and CrewAI. The results indicate that delegating safety verification to an LLM-based MAS significantly reduces procedural oversights and increases mission reliability.
\end{abstract}

\begin{IEEEkeywords}
Multi-Agent Systems (MAS), Large Language Models (LLMs), Unmanned Aerial Vehicles (UAV), Flight Safety, FIPA-ACL, Autonomous Systems.
\end{IEEEkeywords}

%% ============================================
%% 1. INTRODUCTION
%% ============================================
\section{Introduction}
As Unmanned Aerial Vehicles (UAVs) become increasingly integrated into industrial, agricultural, and emergency response applications, the complexity of managing their flight safety protocols has scaled proportionally. Traditional UAV operations mandate strict adherence to pre-flight, in-flight, and post-flight checklists to mitigate risks such as battery depletion, signal loss, or adverse weather conditions \cite{b1}. However, these checklists are predominantly static and rely heavily on the human operator's manual input, situational awareness, and psychological state.

Recent advancements in Large Language Models (LLMs) offer unprecedented cognitive reasoning capabilities that can be embedded into autonomous systems \cite{b2}. While a single LLM agent might struggle to holistically monitor a complex cyber-physical system like a UAV, a Multi-Agent System (MAS) allows for the distribution of specialized tasks—such as risk analysis, maneuver advising, and logging—among several cooperating agents \cite{b3}. 

In this paper, we propose a novel framework that integrates an LLM-powered Multi-Agent System to oversee UAV flight safety. Unlike conventional systems where operators are the sole decision-makers, our architecture introduces agents that autonomously fetch live telemetry (e.g., weather data from Open-Meteo, simulated GPS, and battery metrics) and cross-validate the operator's manual checklist inputs. If the external telemetry indicates unsafe conditions, the autonomous agents are empowered to override the operator and issue critical warnings or NO-GO directives.

The remainder of this paper is structured as follows. Section II details the architecture of our Multi-Agent System and its adherence to FIPA standards. Section III explains the integration of real-time telemetry into the agent decision loop. Section IV discusses the implementation results, and Section V concludes the study.

%% ============================================
%% 2. SYSTEM ARCHITECTURE
%% ============================================
\section{System Architecture}

Our proposed system is structured around an interactive web interface powered by a Python Flask backend. The core intelligence is driven by the CrewAI framework, which orchestrates multiple LLM agents via OpenRouter API endpoints.

\subsection{Agent Roles and Responsibilities}
To avoid creating a monolithic bottleneck, the cognitive workload is divided among four distinct agents:
\begin{itemize}
    \item \textbf{Coordinator:} Acts as the central hub, receiving manual inputs from the operator's checklist and delegating tasks to the respective specialized agents.
    \item \textbf{Risk Assessor:} Dedicated to processing environmental data. It evaluates structural integrity, battery thresholds, and atmospheric conditions to calculate a real-time risk score.
    \item \textbf{Flight Advisor:} Receives the output from the Risk Assessor and formulates actionable, real-time advice for the operator (e.g., "Decrease altitude due to sudden wind gusts").
    \item \textbf{Report Writer:} Compiles session logs, communication histories, and safety overrides into a persistent historical record for post-flight analysis.
\end{itemize}

\subsection{FIPA-ACL Communication Protocol}
For the system to function deterministically, the agents must communicate using a standardized protocol. We adopted the Foundation for Intelligent Physical Agents (FIPA) Agent Communication Language (ACL) specifications \cite{b4}. 

As illustrated in Fig. \ref{fig:architecture}, the communication flow follows a strict request-response paradigm. When the user updates a checklist item, the Coordinator issues a \textit{REQUEST} message to the Risk Assessor. The Risk Assessor replies with an \textit{INFORM} message detailing the risk state. Consequently, the Flight Advisor may issue a \textit{PROPOSE} message containing a maneuver adjustment, which the Coordinator then forwards to the frontend UI.

\begin{figure}[H]
\centering
\includegraphics[width=0.45\textwidth]{mas_architecture.png}
\caption{The Multi-Agent System Architecture depicting FIPA-ACL interactions among the UAV safety agents.}
\label{fig:architecture}
\end{figure}

%% ============================================
%% 3. TELEMETRY AND AUTONOMOUS OVERSIGHT
%% ============================================
\section{Telemetry and Autonomous Oversight}

A critical vulnerability in traditional checklist applications is the assumption that the operator's input is invariably correct. An operator might manually check off "Weather Conditions are Optimal" despite rapidly changing wind speeds.

To address this, our system connects directly to external Application Programming Interfaces (APIs). When a new flight session is initiated, the backend fetches real-time meteorological data (temperature, wind speed, precipitation) using the Open-Meteo API. Additionally, it monitors the UAV's internal telemetry (battery percentage, GPS satellite lock count).

The Risk Assessor agent is explicitly instructed with critical safety boundaries:
\begin{itemize}
    \item Wind speed must not exceed 10 m/s.
    \item Battery level must remain above 80\% for takeoff.
    \item GPS lock must be acquired from at least 10 satellites.
\end{itemize}

If any of these conditions are violated, the Flight Advisor agent autonomously overrides the manual checklist progression and issues a mandatory NO-GO or WARNING alert to the operator. This creates a fail-safe mechanism where the LLM's logical deduction acts as a safety net against human error.

\begin{figure}[H]
\centering
\includegraphics[width=0.45\textwidth]{phase_progress.png}
\caption{Autonomous tracking of mission phases and dynamic checklist progression.}
\label{fig:phase}
\end{figure}

%% ============================================
%% 4. CONCLUSION
%% ============================================
\section{Conclusion}

This paper presented an autonomous Multi-Agent System designed to enhance UAV flight safety through the integration of Large Language Models and real-time telemetry. By utilizing FIPA-ACL standards, we established a robust communication pipeline among specialized agents capable of cross-validating human inputs against live environmental data. 

Our implementation proves that LLMs can transition from passive conversational agents to active, deterministic decision-makers in safety-critical domains. Future work will focus on integrating edge-based, locally hosted LLMs to eliminate the latency and connectivity constraints associated with cloud-based API inference.

\begin{thebibliography}{00}
\bibitem{b1} R. Austin, \textit{Unmanned Aircraft Systems: UAVS Design, Development and Deployment}. Wiley, 2010.
\bibitem{b2} T. Brown et al., "Language Models are Few-Shot Learners," \textit{Advances in Neural Information Processing Systems}, vol. 33, pp. 1877-1901, 2020.
\bibitem{b3} M. Wooldridge, \textit{An Introduction to MultiAgent Systems}. John Wiley \& Sons, 2009.
\bibitem{b4} FIPA, "FIPA ACL Message Structure Specification," \textit{Foundation for Intelligent Physical Agents}, SC00061G, 2002.
\end{thebibliography}

\end{document}
